% !TeX root = ../main.tex
\begin{abstract}
多精度张量计算单元通常采用 $1:2:4$ 的 16、8、4 位峰值算力配比。模型计算向低精度迁移后，增加低精度算力的收益同时取决于硬件成本和工作量组成。本文提出 DP-Puzzle，通过组合共享全精度 Group 与 Q、E、QE 专用 Group 构建九种点积单元配比，覆盖工业界已有选择及更大的增配空间。基于 12\,nm、1\,GHz 下的综合面积和真实负载活动功耗，建立由基准阵列面积占比与功耗预算倍率共同约束的整核部署模型，将单 DP 成本映射为分精度绝对算力。结果表明，$1:2:8$ 的 E 算力翻倍仅增加 33.60\% 的 DP 面积，而 $1:4:8$ 的 Q/E 同时翻倍增加 68.21\%；更大的增配还受到可配置核数的限制。以 DeepSeek-V4.1-Flash 为主要负载，构造 11 个上下文与 32 个算子精度方案，用最坏相对损失评价固定芯片。基准面积占比 30\%、DP 功耗预算为基准 1.5 倍时，原始精度下选择 $1:4:4$，跨上下文最坏损失为 3.91\%；覆盖全部迁移方案时选择 $1:4:8$，最坏损失为 20.05\%。该方法将可组合微架构、双预算成本和精度迁移联系起来，为算力配比选择提供依据。
\end{abstract}
\begin{IEEEkeywords}
多精度算力配比，组合式点积单元，张量计算，面积功耗约束，设计空间探索
\end{IEEEkeywords}
